% =====================================================================
%  PAPER 2 (reframed). The former "make MLMC pay" thesis is inverted by
%  the evidence: both standard accelerations were built and run, and the
%  honest finding is that MLMC does NOT pay for this payoff -- the variance
%  route that pays is CONDITIONING, as single-grid turbocharging, sharpened
%  by exact near-cell (kappa=1) integration. The paper is a positive result
%  (the right method) wrapped around a clean negative one (MLMC's limits).
%
%  Status: ALL sections built; numbers are production-run results. Baseline
%  (3-5) regenerated on the reference machine (seeds 1/2/7/11/23); accelerations
%  (6-7) from production runs (seeds 5/11/23). Compiles with pdflatex (no .bib).
% =====================================================================
\documentclass[11pt]{article}

\usepackage[utf8]{inputenc}
\usepackage[T1]{fontenc}
\usepackage[a4paper,margin=1in]{geometry}
\usepackage{amsmath,amssymb}
\usepackage{booktabs}
\usepackage{graphicx}
\usepackage{xcolor}
\definecolor{linkblue}{rgb}{0.10,0.30,0.55}
\usepackage{hyperref}
\hypersetup{colorlinks=true,allcolors=linkblue}
\setlength{\emergencystretch}{3em}

% ----- baseline numbers (regenerated on the reference machine, seeds 1/2/7/11/23) -----
\newcommand{\VarWe}{0.8532}
\newcommand{\VarWd}{0.8529}
\newcommand{\fwdv}{0.0213}
\newcommand{\fwdn}{0.0175}
\newcommand{\eurMC}{7.9977}
\newcommand{\eurSE}{0.0415}
\newcommand{\eurBS}{7.9656}
\newcommand{\zbs}{0.77}
\newcommand{\telcons}{0.386}
\newcommand{\bA}{0.120}
\newcommand{\bB}{0.219}
\newcommand{\bBmain}{0.202}
\newcommand{\bC}{0.418}
\newcommand{\bD}{0.726}
\newcommand{\gammaval}{1.063}
\newcommand{\costratio}{0.63}
\newcommand{\asianpx}{4.2143}

\title{Turbocharged versus multilevel Monte Carlo for\\
rough-volatility Asian options}

\author{Michael Lumor\thanks{Independent research programme; University of
Salford. ORCID: \href{https://orcid.org/0009-0000-0326-3891}{0009-0000-0326-3891}.
Code: \url{https://github.com/Michaellumor/roughvollab}.}}

\date{Draft --- June 2026}

\begin{document}
\maketitle

\begin{abstract}
Multilevel Monte Carlo (MLMC) reduces the cost of pricing by simulation, but for
rough-volatility models its efficiency is curtailed by the low strong-convergence rate
of the volatility driver. We study arithmetic Asian options in the rough Bergomi model
and ask which acceleration, if any, makes simulation cheaper. We first establish a
baseline: an exact, bias-free multilevel coupling gives a variance-decay rate
$\beta\approx2H$ across the roughness spectrum --- the pathwise bound is tight --- so
naive MLMC does not beat standard Monte Carlo (the cost ratio is $\approx\costratio\times$
at a root-mean-square error of $2.5\times10^{-2}$). We then evaluate, on the same
engine, the two standard accelerations. An \emph{antithetic} coupling leaves the rate
unchanged and is $\approx$10\% more expensive at matched accuracy. A \emph{conditional}
(turbocharging) estimator reduces variance several-fold, but concentrates the reduction
in the single-level estimator rather than the multilevel corrections, so conditional
MLMC does not beat conditional standard Monte Carlo: across tolerances and seeds, plain
conditional standard Monte Carlo at the finest grid is the cheapest of the four
estimators. We further show that integrating the most singular Volterra cell exactly
(the $\kappa=1$ hybrid) sharpens this winner --- because it is bias-limited, a smaller
weak-error constant buys a coarser grid --- for a further $\approx1.3$--$1.5\times$
saving. The recommendation for this payoff is turbocharged standard Monte Carlo on the
$\kappa=1$ variance path; multilevel Monte Carlo does not earn its place. Every result
is validated against analytic limits and regenerates from a single command.
\end{abstract}

\section{Introduction}

The rough-volatility paradigm models log-volatility as a fractional process with a
small Hurst exponent, $H\approx0.1$, reproducing the steep short-maturity
implied-volatility skew and the persistence of realised variance better than classical
Markovian stochastic-volatility models~\cite{gjr2018,bfg2016}. Its defining feature ---
a volatility driver of low H\"older regularity --- is also its computational difficulty:
the model is non-Markovian, has no simple affine structure, and the dominant pricing
tool is therefore Monte Carlo simulation~\cite{bfg2016,mccrickerd2018}.

Monte Carlo cost is the target of the multilevel method of Giles~\cite{giles2008}; 
its efficiency is governed by three rates (see~\cite{giles2015} for a full account):
the weak rate $\alpha$, the variance-decay rate
$\beta$ of the level corrections, and the cost-growth rate $\gamma$. The method attains
the optimal $\mathcal{O}(\varepsilon^{-2})$ complexity when $\beta>\gamma$ and fails to
when $\beta<\gamma$. For diffusions driven by fractional or Gaussian noise the variance
rate is pinned to the strong rate: $\beta$ equals twice the strong convergence
rate~\cite{bfrs2016}. Rough volatility, with its low strong rate, therefore sits in the
regime where MLMC is least helpful --- unless an acceleration can change the picture.

This paper asks, for one demanding payoff --- the arithmetic Asian option, whose value
depends on the time-average of the \emph{underlying} --- which acceleration pays. We
proceed in two movements. The first establishes a clean baseline: with an exact coupling
and no hybrid-scheme correction, $\beta\approx2H$ is tight and naive MLMC is
uncompetitive with standard Monte Carlo, because the arithmetic average preserves rather
than cancels the slowly-decaying Volterra strong error. The second evaluates the two
standard accelerations on the same engine and reports, honestly, where each leads. The
answer is not the one the construction was built to find: neither acceleration makes
MLMC pay. The antithetic coupling does not change the rate; the conditional estimator
reduces variance substantially but does so more in the single-level estimator than in
the multilevel corrections, so the winner is plain conditional standard Monte Carlo at
the finest grid --- turbocharging, outside the multilevel framework. A separate
simulator refinement (exact near-cell integration) then sharpens that winner further.

\paragraph{What is, and is not, new.} That MLMC underperforms for rough volatility is
known, as is the $\beta=2\times(\text{strong rate})$ theory~\cite{bfrs2016}; the baseline
of Sections~\ref{sec:model}--\ref{sec:baseline} confirms this cleanly for a payoff not
previously isolated. Our negative result is therefore not a limitation of the multilevel method but an instance of its own complexity theory: with β<γ, Giles's framework.
β<γ, Giles's framework~\cite{giles2008,giles2015} itself predicts that refinement will not pay, and we confirm this concretely for a payoff where one might have hoped a better coupling would rescue it. Antithetic multilevel estimation~\cite{gilesszpruch2014} and
conditional (``turbocharging'') variance reduction for rough Bergomi~\cite{mccrickerd2018}
are both established. The contribution here is to evaluate both, bias-free, on the
multilevel coupling of this payoff and to report what they actually do: a negative result
for multilevel Monte Carlo --- it does not earn its place even with the best couplings ---
and a positive one identifying the method that does, turbocharged standard Monte Carlo,
which a further refinement makes cheaper still. The value is a clear, reproducible
recommendation for a payoff where the obvious expectation (improve the coupling to rescue
MLMC) turns out to be wrong.

\paragraph{Contributions.}
\begin{enumerate}
\item An exact, bias-free multilevel coupling for rough Bergomi, validated against four
analytic limits (Sections~\ref{sec:model}--\ref{sec:validation}).
\item A clean baseline: $\beta\approx2H$ is tight, naive MLMC does not beat standard
Monte Carlo, and the mechanism is identified (Section~\ref{sec:baseline}).
\item An evaluation of both standard accelerations (Section~\ref{sec:couplings}): the
antithetic coupling leaves $\beta$ unchanged and is net more expensive; the conditional
estimator's variance reduction is concentrated in the single-level estimator, so
conditional MLMC does not beat conditional standard Monte Carlo.
\item The resulting recommendation, with a four-way matched-accuracy cost comparison
(Section~\ref{sec:results}): turbocharged standard Monte Carlo at the finest grid, on the
$\kappa=1$ variance path, is the method of choice; MLMC does not pay for this payoff.
\end{enumerate}

\section{Related work}
\label{sec:related}

\paragraph{Simulation of rough Bergomi.}
Exact Cholesky sampling of the Volterra process is $\mathcal{O}(n^2)$; the hybrid scheme
of Bennedsen, Lunde and Pakkanen~\cite{blp2017} is the standard fast simulator. McCrickerd
and Pakkanen~\cite{mccrickerd2018} add antithetic and conditional (``turbocharging'')
variance reduction for vanilla rough Bergomi pricing --- the two ideas this paper evaluates
in the multilevel, Asian setting; Zhu et al.~\cite{zhu2021} pursue a Markovian
approximation. Our baseline uses the pure $\kappa=0$ Volterra discretisation with a
discrete compensator, which admits an exact multilevel coupling and keeps the
forward-variance identity exact on every grid; Section~\ref{sec:kappa1} additionally adopts
the $\kappa=1$ hybrid for the recommended estimator.

\paragraph{MLMC for rough volatility.}
Bourgey and De Marco~\cite{bourgey2021} apply MLMC to \emph{VIX} options in rough Bergomi
and recover $\mathcal{O}(\varepsilon^{-2})$ with a trapezoidal scheme. Bayer, Ben Hammouda
and Tempone~\cite{bht2020} price vanilla rough Bergomi options and show MLMC performance
``is very sensitive to the values of $H$'' because of low strong-convergence rates, with
\emph{numerical smoothing} restoring the optimal variance decay. Jeng and
Kili\c{c}man~\cite{jk2021} combine MLMC with a control variate for rough Heston. The
general theory is Bayer, Friz, Riedel and Schoenmakers~\cite{bfrs2016}:
$\beta=2\times(\text{strong rate})$ for Gaussian-driven SDEs. Our finding is complementary
to numerical smoothing: rather than restoring MLMC's rate, we find that for the
arithmetic-Asian payoff the single-grid turbocharged estimator is both cheaper than MLMC
and avoids the multilevel machinery.

\paragraph{Positioning.}
The arithmetic-Asian-on-the-underlying payoff is not separately treated in the works above
(their focus is VIX options or European vanillas). We isolate it, show that multilevel
Monte Carlo does not pay for it even with the two standard accelerations, and identify the
method that does.

\section{The rough Bergomi model and an exact multilevel coupling}
\label{sec:model}

\paragraph{Model.}
The rough Bergomi model~\cite{bfg2016} specifies the instantaneous variance through the
Riemann--Liouville (Volterra) process
\begin{equation}
\widetilde{W}_t \;=\; \sqrt{2H}\int_0^t (t-s)^{H-1/2}\,dW_s,
\qquad \operatorname{Var}(\widetilde{W}_t)=t^{2H},
\end{equation}
with Hurst exponent $H\in(0,\tfrac12)$ and standard Brownian motion $W$. The variance and
asset processes are
\begin{equation}
V_t = \xi_0\,\exp\!\Big(\eta\,\widetilde{W}_t-\tfrac12\eta^2 t^{2H}\Big),
\qquad
\frac{dS_t}{S_t} = \sqrt{V_t}\,dB_t,
\quad B=\rho W+\sqrt{1-\rho^2}\,W^{\perp},
\end{equation}
with flat initial forward variance $\xi_0$, volatility-of-volatility $\eta>0$, and
spot--vol correlation $\rho\in[-1,1]$; the compensator makes $\mathbb{E}[V_t]=\xi_0$. We
price the fixed-strike arithmetic Asian call,
$\big(\tfrac1T\!\int_0^T S_t\,dt-K\big)^{+}$.

\paragraph{Discretisation.}
On a uniform grid $t_i=iT/n$ the Volterra process is approximated by the left-point
($\kappa=0$) rule
\begin{equation}
\widetilde{W}^{(n)}_{t_i}
=\sqrt{2H}\sum_{j=0}^{i-1} (t_i-t_j)^{H-1/2}\,\Delta W_j,
\qquad \Delta W_j = W_{t_{j+1}}-W_{t_j},
\label{eq:conv}
\end{equation}
a discrete convolution evaluated in $\mathcal{O}(n\log n)$ by fast Fourier transform.
Compensating with the \emph{discrete} variance
$v_n=\operatorname{Var}(\widetilde{W}^{(n)}_{t_n})$ rather than $t_n^{2H}$ makes
$\mathbb{E}[V^{(n)}_{t_i}]=\xi_0$ hold exactly on every grid.

\paragraph{Exact coupling.}
The multilevel estimator telescopes the price across grids of $n_\ell=n_0 2^{\ell}$ steps,
with level difference $Y_\ell=P^{(\ell)}-P^{(\ell-1)}$ computed on a single Brownian path
shared between the fine grid ($n_\ell$ steps) and the coarse grid ($n_{\ell-1}=n_\ell/2$
steps). The coupling is exact: the coarse increments are the pairwise sums of the fine
increments, $\Delta W^{(c)}_k=\Delta W^{(f)}_{2k}+\Delta W^{(f)}_{2k+1}$, an identity for
Brownian motion. No hybrid-scheme correction enters the coupling, so the level variance
reflects the intrinsic strong error of the scheme and nothing else --- which makes it a
clean reference for the accelerations of Section~\ref{sec:couplings}.

\section{Engine validation}
\label{sec:validation}

The engine is checked against four limits with known answers.
\begin{enumerate}
\item \textbf{Volterra variance.} The empirical variance of $\widetilde{W}^{(n)}_{T}$
matches its discrete analytic value ($\VarWe$ against $\VarWd$; continuum $T^{2H}=1$, the
$\sim15\%$ gap being the expected slow convergence of the rough kernel ---
Section~\ref{sec:kappa1} closes it).
\item \textbf{Forward-variance identity.} With the discrete compensator,
$\max_t|\mathbb{E}[V_t]/\xi_0-1|=\fwdv$, comparable to the Monte-Carlo band
($\approx\fwdn$); the identity $\mathbb{E}[V_t]=\xi_0$ holds exactly by construction, the
small excess being the maximum taken over correlated grid points.
\item \textbf{Black--Scholes limit.} At $\eta=0$ the European price agrees with the closed
form ($\eurMC\pm\eurSE$ against $\eurBS$, $z=\zbs$).
\item \textbf{Telescoping identity.} The normalised consistency check is $\telcons$ in the
main run and stays below $0.51$ in every configuration.
\end{enumerate}

\section{The naive baseline: a tight bound and an uncompetitive cost}
\label{sec:baseline}

\subsection{The variance rate is tight}
The level statistics are estimated from $N$ coupled paths per level, with the rates fitted
by regressing $\log_2\operatorname{Var}(Y_\ell)$, $\log_2|\mathbb{E}[Y_\ell]|$ and
$\log_2 C_\ell$ against $\ell$: $\operatorname{Var}(Y_\ell)\sim 2^{-\beta\ell}$,
$|\mathbb{E}[Y_\ell]|\sim 2^{-\alpha\ell}$, $C_\ell\sim 2^{\gamma\ell}$; for the
$\mathcal{O}(n\log n)$ construction $\gamma\approx\gammaval$. Sweeping $H$
(Table~\ref{tab:betaH}) shows $\beta\approx2H$ throughout: the pathwise bound is tight, the
direct empirical counterpart of $\beta=2\times(\text{strong rate})$~\cite{bfrs2016}. The
$N=20{,}000$ run at $H=0.10$ returns $\beta=\bBmain$, consistent with the sweep. The weak
rate $\alpha$ is noise-dominated at these sample sizes (level biases fall below twice their
standard error) and is not reported as a measurement; the conclusions rest on $\beta$ and
$\gamma$.

\begin{table}[t]
\centering
\caption{Baseline variance-decay rate $\beta$ against the bound $2H$ (production run, seeds
7/11/23).}
\label{tab:betaH}
\begin{tabular}{cccc}
\toprule
$H$ & measured $\beta$ & bound $2H$ & regime \\
\midrule
0.05 & \bA & 0.10 & $\beta<\gamma$ \\
0.10 & \bB & 0.20 & $\beta<\gamma$ \\
0.20 & \bC & 0.40 & $\beta<\gamma$ \\
0.35 & \bD & 0.70 & $\beta<\gamma$ \\
\bottomrule
\end{tabular}
\end{table}

\subsection{Naive MLMC does not pay}
Running the adaptive algorithm of Giles~\cite{giles2008} to a target error $\varepsilon$
and comparing total cost with standard Monte Carlo gives Table~\ref{tab:cost}: at the
tightest tolerance the cost ratio (standard Monte Carlo over MLMC) is $\approx\costratio$
--- naive coupled MLMC costs \emph{more} than plain Monte Carlo. With $\beta=2H$ and
$\gamma\approx1$ the Giles complexity for $\beta<\gamma$ is
$\mathcal{O}(\varepsilon^{-2-(\gamma-\beta)/\alpha})$, strictly worse than
$\mathcal{O}(\varepsilon^{-2})$.

\begin{table}[t]
\centering
\caption{Baseline MLMC cost at fixed convergence rates (production run, seed 11). The finest
level $L$ is monotone in $\varepsilon$, and the cost ratio (standard Monte Carlo over MLMC)
stays below one throughout --- naive MLMC never beats standard Monte Carlo.}
\label{tab:cost}
\begin{tabular}{ccccc}
\toprule
$\varepsilon$ & finest level $L$ & total cost & std MC / MLMC & price \\
\midrule
$2.0\times10^{-1}$ & 2 & $7.06\times10^{5}$ & $0.28$ & $4.1695$ \\
$1.0\times10^{-1}$ & 3 & $3.02\times10^{6}$ & $0.52$ & $4.1996$ \\
$5.0\times10^{-2}$ & 5 & $4.53\times10^{7}$ & $0.55$ & $4.1985$ \\
$2.5\times10^{-2}$ & 7 & $6.32\times10^{8}$ & $0.63$ & $4.2143$ \\
\bottomrule
\end{tabular}
\end{table}

\subsection{The obstacle: the uncancelled Volterra error}
\label{sec:mechanism}
The level difference $Y_\ell$ is driven by
$\widetilde{W}^{(n_\ell)}-\widetilde{W}^{(n_{\ell-1})}$, whose size is the strong error (of
order $n_\ell^{-H}$) and which is strongly correlated \emph{along the path}: the same
coarse-grid misrepresentation of the volatility trajectory shifts $S_t$ in the same
direction at nearby times. The arithmetic average $\tfrac1T\int_0^T S_t\,dt$ therefore does
not average the error away; it preserves the common component, so
$\operatorname{Var}(Y_\ell)$ inherits the single-point rate $2H$. This is the obstacle the
accelerations target: each acts on the integrand \emph{before} averaging, with the aim of
cancelling or removing the common error.

\section{Evaluating two accelerations}
\label{sec:couplings}

Both accelerations preserve the bias-free property ($\mathbb{E}[V_t]=\xi_0$ on every grid,
exact increment summation across levels), so the telescoping identity of
Section~\ref{sec:validation} still holds and each estimator is unbiased by construction. We
measure each against the baseline, and --- crucially --- compare costs at matched accuracy
(the methodological reason is in Section~\ref{sec:antithetic}).

\subsection{Antithetic coupling}
\label{sec:antithetic}
Following the antithetic multilevel construction of Giles and
Szpruch~\cite{gilesszpruch2014}, the coarse-level estimator is replaced by the average of
two coupled fine paths: the standard fine path and an antithetic path in which the two fine
Brownian increments within each coarse step are swapped,
$(\Delta W^{(f)}_{2k},\Delta W^{(f)}_{2k+1})\mapsto(\Delta W^{(f)}_{2k+1},\Delta
W^{(f)}_{2k})$ for both drivers. The coarse increment is unchanged (the sum is invariant),
so the coupling stays exact. For a classical diffusion this cancels the leading strong
error and lifts $\beta$; for the rough Volterra process it does not. The swap leaves the
leading linear kernel error symmetric, so it is not cancelled, and the variance rate is
unchanged: across $H$, $\beta_{\text{anti}}$ equals the naive $\beta$ to two decimals,
tracking $2H$ and nowhere near $4H$ (Table~\ref{tab:betaH}, each row reproduced by the
antithetic estimator). The variance is reduced by a constant factor only --- a mean of
$1.44\times$ (range $1.37$--$1.49\times$ over $H$). At matched accuracy this does not pay:
the antithetic path adds a second fine-grid evaluation per coupled level, and the
$\approx1.44\times$ variance reduction falls short of that cost increase, leaving the
estimator $\approx9$--$11\%$ more expensive (efficiency $\approx0.91\times$ at the matched
finest level, across tolerances).

\paragraph{A note on cost comparison.} Comparing two estimators with identical bias requires
holding the finest level fixed. An adaptive selector chooses the finest level from the
measured variances, and because the antithetic estimator has lower per-level variance it
stops at a coarser level than the naive one --- manufacturing an apparent advantage that has
nothing to do with efficiency at matched accuracy. We observed this directly: the
free-running antithetic/naive cost ratio swung from $0.13\times$ to $1.10\times$ across
tolerances purely through level selection (the estimators choosing $L=2$ versus $L=5$ at one
tolerance), while the bias at the coarser level was measurably larger. The effect flips sign
across random seeds, confirming it is an artifact. We therefore pin the finest level for all
estimator comparisons, here and in Section~\ref{sec:results}; the cost figures are
matched-accuracy throughout.

% ===================== conditional subsection (cited results) =====================
\subsection{Conditional Monte Carlo}
\label{sec:conditional}

The second acceleration conditions on the variance path to remove the orthogonal driving
noise. Given the path of $W$ (hence of $V$), the log-asset is Gaussian, so the geometric
average $G=\exp\!\big(\sum_i w_i\log S_{t_i}\big)$ is conditionally lognormal and the
geometric-Asian value $\mathbb{E}[(G-K)^+\mid W]$ is available in closed form. We use it as
a control variate with unit coefficient,
\begin{equation}
P_{\mathrm{cond}} \;=\; (\bar A-K)^+ \;-\;
\big[(G-K)^+ - \mathbb{E}[(G-K)^+\mid W]\big],
\label{eq:condcv}
\end{equation}
applied on each grid; $\bar A$ and $G$ share the trapezoidal weights of the engine, and the
conditional mean and variance of $\log G$ are deterministic functions of the simulated
variance path. The estimator is unbiased by construction, and we confirm it reproduces the
baseline price to Monte-Carlo precision.

Conditioning removes a large but \emph{non-uniform} share of the variance. Paired and
seed-averaged, it reduces the single-level payoff variance by $4.16\times$ (mean over
levels; $4.21\times$ at the finest grids) but the level-difference variance
$\operatorname{Var}(Y_\ell)$ by only $3.19\times$ ($3.30\times$ at the finest). The
single-level reduction holds flat at $4.1$--$4.25\times$ across $\ell=0,\dots,6$, while the
level-difference reduction climbs from $2.9$ to $3.3\times$. The single-level estimator is
reduced more than the level difference because the orthogonal-driver variance that
conditioning eliminates is precisely the part the level difference has \emph{already} partly
cancelled; what remains in $Y_\ell$ still decays at the rough rate $\beta\approx2H\ll\gamma$.
Conditioning lowers the constant but does not touch the pathology, so multilevel still cannot
repay refinement.

The cost comparison makes this concrete. We pin the finest level $L^{\*}$ at each tolerance
and compare the optimal multilevel cost
$(2/\varepsilon^2)\big(\sum_\ell\sqrt{V_\ell C_\ell}\big)^2$ against the single-level cost
$(2/\varepsilon^2)\operatorname{Var}(P_{L^{\*}})\,C_{L^{\*}}$, each with and without
conditioning. With a measured per-path conditioning overhead $\kappa=1.275$, the
matched-accuracy work (units of fine-step evaluations, seed-averaged) is:

\begin{table}[h]
\centering
\caption{Matched-accuracy cost of the four estimators (pinned $L^{\*}$, seed-averaged,
$\kappa=1.275$). Conditional standard Monte Carlo is cheapest at both tolerances and on every
seed.}
\label{tab:fourway}
\begin{tabular}{cccccc}
\toprule
$\varepsilon$ & $L^{\*}$ & naive ML & cond.\ ML & naive std & cond.\ std \\
\midrule
$0.10$ & 2 & $1.40\times10^{6}$ & $5.31\times10^{5}$ & $7.81\times10^{5}$ & $\mathbf{2.38\times10^{5}}$ \\
$0.05$ & 5 & $4.64\times10^{7}$ & $1.80\times10^{7}$ & $2.46\times10^{7}$ & $\mathbf{7.50\times10^{6}}$ \\
\bottomrule
\end{tabular}
\end{table}

At both tolerances and across every seed the ordering is identical: conditional standard
Monte Carlo is the cheapest of the four, conditional multilevel next, then naive standard,
then naive multilevel. Conditional standard MC is about $3.3\times$ cheaper than naive
standard MC --- the factor $\kappa/4.16\approx0.30$ --- and several times cheaper than any
multilevel variant. Most pointedly, the ratio of conditional-standard to
conditional-multilevel cost is $0.41$--$0.45$ across all seeds: \emph{conditional multilevel
Monte Carlo does not beat conditional standard Monte Carlo}. This ratio is invariant to
$\kappa$, since both estimators pay the same per-path overhead, so the conclusion is robust
to how the conditioning work is costed.

The verdict for this payoff is therefore unambiguous. Of the two variance routes, the one
that pays is conditioning, not multilevel; and once one conditions, the residual
level-difference variance decays too slowly ($\beta\approx2H$) for refinement to be
worthwhile. Conditional (turbocharged) standard Monte Carlo at the finest grid is the method
of choice, and multilevel Monte Carlo does not earn its place.

% ===================== kappa=1 refinement (cited results) =====================
\subsection{Exact near-cell integration sharpens the recommended estimator}
\label{sec:kappa1}

The recommendation above can be sharpened by improving the simulator rather than the
estimator. The hybrid scheme used so far approximates the whole Volterra kernel by an
optimally weighted Riemann sum; integrating the nearest, most singular cell exactly instead
(the $\kappa=1$ hybrid) closes essentially all of the discrete variance error --- the
simulated $\operatorname{Var}(\widetilde W_T)$ rises from $0.853$ to $0.9996$ against the
continuum $T^{2H}=1$, the near-cell singularity having accounted for the entire gap.
Crucially this is a smaller \emph{bias} constant, not a faster rate: the strong rate, and
with it $\beta\approx2H$, are unchanged.

A smaller bias constant helps precisely when the estimator is bias-limited, and conditional
standard Monte Carlo is. Conditioning removes enough sampling variance (\S\ref{sec:conditional})
that the finest grid is set by the weak-error requirement, not the sample budget: at target
RMSE $\varepsilon$ the bias-determined grid is $n^{\*}=32$ ($\varepsilon=0.05$) and
$n^{\*}=64$ ($\varepsilon=0.025$), and the next-coarser grid violates the bias budget
regardless of how many paths are drawn. The $\kappa=1$ scheme cuts the discretisation bias to
$\approx0.59\times$, which clears the same budget one grid coarser --- $n^{\*}\!:32\to16$ and
$64\to32$ --- at a cost of only a $\approx1.13\times$ larger single-level variance and a
measured $1.03$--$1.12\times$ per-path overhead (the extra per-cell Gaussian is cheap beside
the FFT). The grid halving more than repays both:

\begin{table}[h]
\centering
\caption{Matched-accuracy cost of conditional standard MC on the $\kappa=0$ vs $\kappa=1$
variance path (bias-determined grid $n^{\*}$, seed-averaged). The $\kappa=1$ saving widens as
$\varepsilon$ tightens.}
\label{tab:kappa1cost}
\begin{tabular}{cccc}
\toprule
$\varepsilon$ & $n^{\*}$ ($\kappa{=}0$) & $n^{\*}$ ($\kappa{=}1$) & cost ratio $\kappa{=}1/\kappa{=}0$ \\
\midrule
$0.05$  & 32 & 16 & $0.79$ \\
$0.025$ & 64 & 32 & $0.68$ \\
\bottomrule
\end{tabular}
\end{table}

The matched-accuracy cost ratio is $0.79$ at $\varepsilon=0.05$ and $0.68$ at
$\varepsilon=0.025$ --- a $1.3$--$1.5\times$ saving that widens as the tolerance tightens,
and stable across seeds. We stress the scope: the $\kappa=1$ refinement does \emph{not} help
multilevel Monte Carlo, where it raises the per-level variance at unchanged rate and is
strictly unfavourable; its value is confined to the single-grid estimator that the
rough-volatility cost structure already selected. It is a constant-factor improvement to the
method of choice, not a change to the ranking of methods.

\section{Synthesis and recommendation}
\label{sec:results}

The three evaluations converge on a single picture, governed by one rate. Naive multilevel
Monte Carlo loses because $\beta\approx2H<\gamma$ (Section~\ref{sec:baseline}). The antithetic
coupling does not change $\beta$ and is net more expensive (Section~\ref{sec:antithetic}). The
conditional estimator reduces variance several-fold, but concentrates the reduction in the
single-level estimator rather than the multilevel corrections, so conditional MLMC does not
beat conditional standard Monte Carlo (Section~\ref{sec:conditional}). Exact near-cell
integration then sharpens the single-grid winner (Section~\ref{sec:kappa1}). The four-way
comparison of Table~\ref{tab:fourway} is decisive at every tolerance and seed: conditional
standard Monte Carlo at the finest grid is the cheapest estimator, and the $\kappa=1$ variance
path makes it cheaper still.

The recommendation for arithmetic Asian options under rough Bergomi is therefore turbocharged
standard Monte Carlo on the $\kappa=1$ variance path. Multilevel Monte Carlo does not earn its
place for this payoff --- not because the coupling is poor (it is exact and tight) but because
the rate $\beta\approx2H$ is intrinsic to the rough kernel and neither acceleration changes it.
This is consistent with, and complementary to, the numerical-smoothing route of
Bayer--Ben~Hammouda--Tempone~\cite{bht2020}, which restores MLMC's rate by smoothing the
integrand; we find that for this payoff the simpler single-grid turbocharged estimator is both
cheaper and avoids the multilevel machinery entirely.

\section{Discussion}
\label{sec:discussion}

\paragraph{What this establishes.} For arithmetic Asian options in rough Bergomi, the pathwise
variance bound $\beta\approx2H$ is tight, and neither of the two standard accelerations lifts
it. The variance reduction that pays --- conditioning --- benefits the single-grid estimator
more than the multilevel one, because the orthogonal noise it removes is concentrated in the
single-level variance. The practical consequence is a clear recommendation (turbocharged
standard Monte Carlo, on the $\kappa=1$ variance path) and a clean negative result (multilevel
Monte Carlo does not pay here), both measured at matched accuracy and stable across seeds.

\paragraph{Limitations.} The weak rate $\alpha$ is noise-dominated and not reported as a
measurement; the study fixes a single regime, payoff, and coarsest grid. The cost comparisons
are at matched accuracy with the finest level pinned, which we argue is the only honest basis
for comparing equal-bias estimators; the free-running adaptive figures are reported only to
illustrate the artifact. The FFT-based construction makes the numbers reproducible to
Monte-Carlo precision with documented seeds, not bitwise across environments.

\section{Conclusion}
\label{sec:conclusion}

We set out to determine whether multilevel Monte Carlo could be made to pay for the arithmetic
Asian option in rough Bergomi, and found, on the evidence of both standard accelerations built
and run on an exact coupling, that it cannot: the variance rate $\beta\approx2H$ is intrinsic
and neither antithetic nor conditional coupling changes it. The acceleration that does pay is
conditioning, but as single-grid turbocharging rather than multilevel --- conditional standard
Monte Carlo at the finest grid is the cheapest estimator across tolerances and seeds, and
exact near-cell integration sharpens it by a further $1.3$--$1.5\times$. The recommendation is
specific and reproducible, and the negative result for multilevel Monte Carlo is, for this
payoff, as informative as the positive one.

\section*{Reproducibility statement}

All numerical results are produced by \texttt{layer1b\_mlmc\_asian.py} in the public
repository (\url{https://github.com/Michaellumor/roughvollab}), authored by the named author
(ORCID 0009-0000-0326-3891). The simulator is unit-tested against the analytic limits of
Section~\ref{sec:validation}; sections run from a single entry point with documented random
seeds (7/11/23 for the baseline and antithetic gates, 5/11/23 for the conditional and
$\kappa=1$ comparisons). Because the Volterra process is built by FFT convolution, results
reproduce to Monte-Carlo precision with those seeds rather than bitwise across environments;
the adaptive selector is additionally noise-sensitive at intermediate tolerances, which is why
all estimator comparisons pin the finest level. Baseline numbers are to be regenerated on the
author's reference machine; the acceleration results are from production runs on that engine.

\section*{Use of generative AI}

This work was produced with substantial assistance from an AI assistant (Anthropic's Claude),
used for writing and refactoring code, drafting and editing prose, and suggesting analyses and
checks. The research direction, the modelling and implementation decisions, the interpretation
of results, and the responsibility for their correctness are the author's own.

\begin{thebibliography}{99}

\bibitem{bht2020} Bayer, C., Ben Hammouda, C., Tempone, R. (2020). Hierarchical adaptive
sparse grids and quasi-Monte Carlo for option pricing under the rough Bergomi model.
\textit{Quantitative Finance} 20(9), 1457--1473.

\bibitem{bfg2016} Bayer, C., Friz, P., Gatheral, J. (2016). Pricing under rough volatility.
\textit{Quantitative Finance} 16(6), 887--904.

\bibitem{bfrs2016} Bayer, C., Friz, P.~K., Riedel, S., Schoenmakers, J. (2016). From rough
path estimates to multilevel Monte Carlo. \textit{SIAM Journal on Numerical Analysis} 54(3),
1449--1483.

\bibitem{blp2017} Bennedsen, M., Lunde, A., Pakkanen, M.~S. (2017). Hybrid scheme for
Brownian semistationary processes. \textit{Finance and Stochastics} 21(4), 931--965.

\bibitem{bourgey2021} Bourgey, F., De Marco, S. (2021). Multilevel Monte Carlo simulation for
VIX options in the rough Bergomi model. \textit{Journal of Computational Finance};
arXiv:2105.05356.

\bibitem{gjr2018} Gatheral, J., Jaisson, T., Rosenbaum, M. (2018). Volatility is rough.
\textit{Quantitative Finance} 18(6), 933--949.

\bibitem{giles2008} Giles, M.~B. (2008). Multilevel Monte Carlo path simulation.
\textit{Operations Research} 56(3), 607--617.

\bibitem{giles2015} Giles, M.~B. (2015). Multilevel Monte Carlo methods.
\textit{Acta Numerica} 24, 259--328.

\bibitem{gilesszpruch2014} Giles, M.~B., Szpruch, L. (2014). Antithetic multilevel Monte
Carlo estimation for multidimensional SDEs without L\'evy area simulation. \textit{Annals of
Applied Probability} 24(4), 1585--1620.

\bibitem{jk2021} Jeng, S.~W., Kili\c{c}man, A. (2021). On multilevel and control variate
Monte Carlo methods for option pricing under the rough Heston model. \textit{Mathematics}
9(22), 2930.

\bibitem{mccrickerd2018} McCrickerd, R., Pakkanen, M.~S. (2018). Turbocharging Monte Carlo
pricing for the rough Bergomi model. \textit{Quantitative Finance} 18(11), 1877--1886;
arXiv:1708.02563.

\bibitem{zhu2021} Zhu, Q., Loeper, G., Chen, W., Langren\'e, N. (2021). Markovian
approximation of the rough Bergomi model for Monte Carlo option pricing. \textit{Mathematics}
9(5), 528.

\end{thebibliography}

\end{document}
